\section{Sparse Readout Features Capture Nontrivial Structure}
\label{sec:results}
\label{subsec:calibration}

This section builds one cumulative argument: selected-score
reconstruction is accurate in a declared operating regime; simpler
geometry recovers neither the same score fidelity nor the same
groupings; the explanation-level structure reproduces across independent
training runs; and feature contributions predict realized local
intervention effects. Feature decompositions are read only when
reconstructed and exact signed scores agree and satisfy \(\rho<0.5\)
(\Cref{sec:method-reading-local-accounts}).

\paragraph{Readout and Score Fidelity}\label{subsec:readout-replacement-fidelity}
\begin{table}[t]
\caption{\textbf{Readout replacement and selected-score reconstruction.}
Q/G/Min/R1-Q/R1-L \(=\) Qwen3.5 / Gemma-4 / Ministral-3 /
R1-Distill-Qwen / R1-Distill-Llama (suffix is parameter count).
cfg: width multiplier / TopK budget. rowEV (row-centered explained variance):
fraction of variance in row-centered unembedding rows captured by the sparse
reconstruction. top1: fraction of held-out hidden states where SRP and the
dense LM head agree on the argmax token. KL: median KL divergence between
dense and SRP-reconstructed next-token distributions. Sign: fraction of
contrasts where the SRP-reconstructed margin shares sign with the exact
margin. Replacement uses held-out hidden states; score metrics use the main
held-out logit-difference set at each model's \(k=256\) operating point.
Med.\ \(\rho\) and \(\rho<0.5\) use the floored \(\rho\) of
\S\ref{sec:method-reading-local-accounts}; cluster-bootstrap 95\% CIs are
in \Cref{tab:app-fidelity-cis}. Models above the rule pass the
interpretation gate (top1 \(\ge 0.75\)); the two Gemma variants below it
do not. Higher is better for rowEV, top1, Sign, and
\(\rho<0.5\); lower is better for KL and median \(\rho\). Full names:
Appendix~\ref{app:model-suite-selection}.}
\label{tab:readout-score-fidelity-summary}
\centering
\scriptsize
\setlength{\tabcolsep}{2.0pt}
\renewcommand{\arraystretch}{1.06}
\begin{tabular}{@{}lc
S[table-format=1.3]
S[table-format=1.3]
S[table-format=1.3]
S[table-format=1.3]
S[table-format=1.3]
S[table-format=1.3]@{}}
\toprule
& & \multicolumn{3}{c}{Replacement} & \multicolumn{3}{c}{Logit differences} \\
\cmidrule(lr){3-5}\cmidrule(l){6-8}
Model & cfg & {rowEV} & {top1} & {KL} & {Sign} & {Med. \(\rho\)} & {\(\rho<0.5\)} \\
\midrule
Q-0.8B & 32/256 & 0.877 & 0.891 & 0.135 & 0.950 & 0.075 & 0.891 \\
Q-2B   & 32/256 & 0.847 & 0.887 & 0.136 & 0.958 & 0.071 & 0.890 \\
Q-9B   & 16/256 & 0.857 & 0.900 & 0.105 & 0.949 & 0.088 & 0.878 \\
Min-8B & 32/256 & 0.888 & 0.904 & 0.087 & 0.949 & 0.068 & 0.888 \\
R1-Q-7B & 32/256 & 0.844 & 0.760 & 0.489 & 0.921 & 0.147 & 0.787 \\
R1-L-8B & 32/256 & 0.888 & 0.754 & 0.434 & 0.926 & 0.134 & 0.830 \\
\midrule
G-E2B  & 32/256 & 0.834 & 0.333 & 6.37  & 0.914 & 0.255 & 0.710 \\
G-E4B  & 32/256 & 0.827 & 0.736 & 1.22  & 0.940 & 0.155 & 0.807 \\
\bottomrule
\end{tabular}
\end{table}

\label{subsec:selected-score-reconstruction}%
The factorization preserves both the readout and the selected scores
(\Cref{tab:readout-score-fidelity-summary}). Replacement fidelity, not
aggregate row capacity, licenses interpretation: each readout is certified
by top-token agreement and KL on held-out hidden states before any local
decomposition, since rows with large residual norms can dominate the
softmax even when reconstruction is accurate in aggregate
(Appendix~\ref{app:readout-norm-tails}). At the score level, sign
agreement is at least \(0.91\) across the suite and the low-error fraction
(\(\rho<0.5\)) covers \(0.71\)--\(0.89\) of contrasts per model, with
residual error and the few sign flips confined to near-tie margins
(Appendix Figure~\ref{fig:readout-score-fidelity-result}); the case
studies in later sections are drawn from the passing subset, and
small-margin scores remain in the aggregate
statistics.\label{subsec:score-family-fidelity} Fidelity is predictable
from the score family---difficulty tracks how close the competitor sits
to the selected token, from perfect sign agreement on vocabulary-mean
contrasts down to \(0.824\) on near-tie top1-vs-top2 margins---because
the contrast residual \(h^\top(r_A-r_B)\) contains no term that scales
with the margin, so near ties are read under the absolute bound of
\S\ref{sec:method-reading-local-accounts}
(Appendix Table~\ref{tab:readout-score-families}). The same fidelity
extends beyond hand-picked token pairs: 300 benchmark-derived contrasts
from SQuAD2, HotpotQA, LegalBench (ContractNLI), and SecurityEval
\citep{rajpurkar2018squad2,yang2018hotpotqa,koreeda2021contractnli,guha2023legalbench,siddiq2022securityeval}
reach median \(\tilde\rho=0.129\) at sign agreement \(0.943\) on Qwen3.5-2B
(Appendix~\ref{app:benchmark-derived-display-examples},
Table~\ref{tab:benchmark-derived-readout-contrasts}).

\paragraph{Simpler Geometry Does Not Recover the Scores}\label{subsec:baseline-comparisons}
\begin{table}[t]
\caption{\textbf{Direct-geometry alternatives across the six gated
models.} Sign-preserved low-error coverage (sign agreement and unfloored
\(\tilde\rho<0.5\); \S\ref{sec:method-reading-local-accounts}) on the
full logit-difference bank (\(\sim\)1{,}350 items per model), with
cluster-bootstrap 95\% CIs. ``Best alternative'' is the strongest, per
model, of six methods: nearest-row ridge, weighted row-kNN, k-means
centroid dictionaries at \(D{=}65{,}536\) and \(D{=}16{,}384\) (both
\(k{=}256\), sparsity-matched), hard cluster assignment, and PCA-256;
full grids in Appendix~\ref{app:direct-geometry-grid}. The best
alternative is nearest-row ridge everywhere except R1-L-8B (k-means
centroid dictionary, \(D{=}65{,}536\)). Lead in percentage points.}
\label{tab:main-baseline-null-comparison}
\centering
\scriptsize
\setlength{\tabcolsep}{2.6pt}
\renewcommand{\arraystretch}{1.06}
\begin{tabular}{@{}lccc@{}}
\toprule
Model & SRP & Best alternative & Lead \\
\midrule
Q-0.8B & 0.754 [.733, .773] & 0.665 [.645, .686] & +8.9 \\
Q-2B   & 0.780 [.759, .798] & 0.625 [.606, .644] & +15.5 \\
Q-9B   & 0.812 [.792, .831] & 0.639 [.618, .658] & +17.3 \\
Min-8B & 0.777 [.759, .795] & 0.648 [.629, .665] & +12.9 \\
R1-Q-7B & 0.716 [.694, .738] & 0.579 [.560, .598] & +13.7 \\
R1-L-8B & 0.751 [.728, .772] & 0.624 [.606, .642] & +12.7 \\
\bottomrule
\end{tabular}
\end{table}

If SRP's codes merely repackaged \(\WU\) geometry, methods built directly
from that geometry should reconstruct selected scores equally well. They
do not (\Cref{tab:main-baseline-null-comparison}): across the six gated
models, SRP's sign-preserved low-error coverage is 0.72--0.81, leading
the strongest of six alternatives by 8.9--17.3 points with
non-overlapping intervals. The comparison is capacity-fair: the k-means
centroid dictionary at \(D{=}65{,}536\), \(k{=}256\) matches SRP's width
and sparsity exactly and still trails by 18 points on the display model,
and shuffled-code and random-pattern nulls collapse entirely
(\Cref{tab:app-cross-model-nulls}). The gap holds at the grouping level
as well: on the cores of feature-group structure that reproduce across
training seeds, a held-out same-recipe SRP dictionary recovers 72--75\%
of the core, against 43--49\% for row-kNN and 21--25\% for hard
clustering, with ratio CIs excluding parity
(Appendix~\ref{app:basis-stability}). Learned sparse row--code structure,
not generic row proximity, carries the reconstruction.

\paragraph{Explanations Reproduce Across Independent Runs}\label{subsec:stability}
Individual decoder directions are not canonical: across three
independently trained dictionaries per width (equal held-out
reconstruction verified), matched decoder cosines for used features
average only \(\sim\)0.33 at the 32\(\times\) width, consistent with the
reported seed-sensitivity of TopK dictionaries
\citep{paulo2025stability}. The interpretive unit, however, reproduces.
Over 63 curated contrasts, the side token-sets of the decomposition
reproduce across every seed pair at \(\sim\)15\(\times\) the
cross-contrast null (100\% of contrasts above the null's 90th percentile
at both widths, cross-side leakage \(\le 0.005\)), and 87--91\% of used
feature groups have an above-chance counterpart in an independently
trained dictionary, the median group 83--90\% covered by a union of at
most three features. Structure reproduces, indexing does not, and the
residual instability is feature splitting rather than noise. These
statements are within-recipe (fixed width, sparsity, and preprocessing);
cross-recipe variation is larger and is disclosed in
Appendix~\ref{app:basis-stability}.

\paragraph{Contributions Predict Local Intervention Effects}\label{subsec:causal-validation}
\begin{table}[t]
\caption{\textbf{Predicted versus realized local readout-side changes.}
For each contrast, the predicted per-feature contribution
\(c_i=\beta_i(q)\,p_i(h)\) is compared with the realized margin change
when \(h\) is ablated along the unit decoder direction \(d_i\)
(readout-side, linear in \(h\); no re-forward). Per model:
\(\sim\)260 bank contrasts, top-10 features plus 10 random-direction
controls each; \(r^2\) over reading-rule-gated pairs with bootstrap
95\% CIs; slope from the same fit; Random: \(r^2\) for the
random-direction controls. Protocol:
Appendix~\ref{app:causal-validation}.}
\label{tab:causal-validation-summary}
\centering
\scriptsize
\setlength{\tabcolsep}{3.2pt}
\renewcommand{\arraystretch}{1.06}
\begin{tabular}{@{}lcccc@{}}
\toprule
Model & \(r^2\) [95\% CI] & Slope & Random & \(n\) \\
\midrule
Q-0.8B & 0.886 [.869, .903] & 1.661 & 0.010 & 1610 \\
Q-2B   & 0.834 [.802, .862] & 1.338 & 0.007 & 1820 \\
Q-9B   & 0.829 [.813, .845] & 1.686 & 0.058 & 1860 \\
Min-8B & 0.926 [.913, .937] & 0.957 & 0.009 & 2050 \\
R1-Q-7B & 0.918 [.899, .933] & 1.082 & 0.000 & 1500 \\
R1-L-8B & 0.934 [.920, .947] & 1.126 & 0.034 & 1210 \\
\bottomrule
\end{tabular}
\end{table}

The decomposition is not merely descriptive bookkeeping: on all six gated
models, predicted feature contributions match the realized margin change
under readout-side ablation at \(r^2=0.83\)--\(0.93\), while random
directions explain at most 0.06
(\Cref{tab:causal-validation-summary}). Sparse-code misattribution,
decoder non-orthogonality, and the row residual would each break this
agreement; none does. Magnitudes are calibrated (slope \(\approx\)1) on
the 7--8B models and overstated by up to \(\sim\)1.7\(\times\) on the
small Qwen readouts with sign and rank preserved---consistent with
greater feature interference in narrower readouts; we report slopes as
measured and leave the diagnosis open. Gating changes \(r^2\) by less
than 0.02. The scope is deliberate: these are local, readout-side
predictions on realized forward passes, not circuit-level necessity or
generation-level effects.

\paragraph{Feature Labels Survive a Blinded Audit}\label{subsec:label-audit}
Feature labels are readability summaries of top unembedding rows, audited
separately from the measured feature ids and signed terms. Three
independent blinded rater runs of a frontier language model (frozen
prompt, top-20 rows only, fresh contexts) agree unanimously on 93.4\% of
items (Fleiss' \(\kappa=0.87\)); by consensus, 76 of the 99 claim-bearing
display labels name a coherent lexical family, 4 are mixed or ambiguous,
and 19 are token-form surface artifacts, in 81.8\% agreement with an
independent author audit
(Appendix~\ref{app:main-case-study-feature-audit}). We describe this as a
same-model reproducibility audit rather than human annotation; claims are
tied to feature ids, signed contributions, and residuals, with labels
used only where audited
\citep{karvonen2025saebench,makelov2024principled}.

\paragraph{The Operating Regime for Term-Level Reading}\label{subsec:interpretation-boundary}
A faithful decomposition is a local reading of the readout side: a
different factorizer or active budget could split the same contribution
across other features, so the fidelity checks tie the reading to
this factorizer and active budget. Appendix~\ref{app:robustness-operating-regime}
routes each gated condition (residual-dominated, sign-flipped,
model-unsupported, diffuse, or token-confounded scores) to the part that
still holds; where the residual dominates a decomposition or its terms
repeatedly change sign, the score is reported as uninterpretable rather
than silently dropped. These conditions bind most often where replacement
fidelity is lowest (\Cref{tab:readout-score-fidelity-summary}); the
Qwen3.5-2B readout used for the displays clears the interpretation checks
by a wide margin.

\paragraph{Further Uses of the Basis}\label{sec:feature-resolved-component-decomposition}\label{subsec:constrained-readout-edits}
Two supporting applications live in appendices. Because signed readout
terms are linear in \(h\), direct logit attribution composes with the
fixed basis, refining each residual-stream component's contribution into
per-feature terms; Appendix~\ref{app:feature-resolved-dla} states the
identity and gives worked displays. The learned decoder directions also
serve as constrained edit directions: on Qwen3.5-2B, dictionary
directions suppress held-out lexical variants more strongly than mean-row
and PCA control directions at every matched level of next-token
distribution shift (12 of 12 paired comparisons significant), but the
advantage is not universal---the 0.8B and 9B Qwen readouts lose to the
mean-row direction, and two candidate explanations (replacement fidelity
and dictionary width) were tested and rejected. Protocol, frontiers, all
cross-model losses, and per-tokenizer exclusions are in
Appendix~\ref{app:lexical-control-stress-test}.

\FloatBarrier
